\documentclass[12pt,a4paper]{article}

%==============================================================
%  PACKAGE
%==============================================================
\usepackage[utf8]{inputenc}
\usepackage[T1]{fontenc}
\usepackage{amsmath, amssymb, amsthm}
\usepackage{geometry}
\usepackage{hyperref}
\usepackage{listings}
\usepackage{xcolor}
\usepackage{booktabs}
\usepackage{array}
\usepackage{tcolorbox}
\usepackage{titlesec}
\usepackage{fancyhdr}

\geometry{top=2.5cm, bottom=2.5cm, left=3cm, right=3cm}

%==============================================================
%  COULEURS
%==============================================================
\definecolor{deepblue}{RGB}{0,51,102}
\definecolor{cblue}{RGB}{0,90,160}
\definecolor{codegreen}{rgb}{0,0.6,0}
\definecolor{codegray}{rgb}{0.5,0.5,0.5}
\definecolor{codepurple}{rgb}{0.58,0,0.82}
\definecolor{backcolour}{rgb}{0.97,0.97,0.97}

%==============================================================
%  CODE STYLE
%==============================================================
\lstdefinestyle{pythonstyle}{
    backgroundcolor=\color{backcolour},
    commentstyle=\color{codegreen}\itshape,
    keywordstyle=\color{cblue}\bfseries,
    numberstyle=\tiny\color{codegray},
    stringstyle=\color{codepurple},
    basicstyle=\ttfamily\small,
    breaklines=true,
    breakatwhitespace=false,
    captionpos=b,
    numbers=left,
    numbersep=5pt,
    frame=single,
    rulecolor=\color{gray!40},
    language=Python,
    showstringspaces=false,
    literate=
      {à}{{\`a}}1 {é}{{\'e}}1 {è}{{\`e}}1 {ê}{{\^e}}1
      {î}{{\^i}}1 {ô}{{\^o}}1 {û}{{\^u}}1 {ù}{{\`u}}1
      {ç}{{\c{c}}}1 {É}{{\'E}}1 {È}{{\`E}}1 {À}{{\`A}}1
}
\lstset{style=pythonstyle}

%==============================================================
%  THEOREM ENVIRONMENTS
%==============================================================
\newtheorem{theoreme}{Théorème}[section]
\newtheorem{corollaire}[theoreme]{Corollaire}
\newtheorem{proposition}[theoreme]{Proposition}
\theoremstyle{definition}
\newtheorem{definition}[theoreme]{Définition}
\newtheorem{exemple}[theoreme]{Exemple}
\newtheorem{hypothese}{Hypothèse}[section]
\theoremstyle{remark}
\newtheorem{remarque}[theoreme]{Remarque}

\newenvironment{demonstration}[1][]
  {\begin{proof}[\textbf{Démonstration}\ifx&#1&\else\ (#1)\fi]}
  {\end{proof}}

%==============================================================
%  SECTION FORMATTING
%==============================================================
\titleformat{\section}{\large\bfseries\color{deepblue}}{}{0em}{\thesection\quad}
\titleformat{\subsection}{\normalsize\bfseries\color{cblue}}{}{0em}{\thesubsection\quad}

%==============================================================
%  HEADER FOOTER
%==============================================================
\pagestyle{fancy}
\fancyhf{}
\rhead{\small\textcolor{gray}{La Régression Linéaire Simple}}
\lhead{\small\textcolor{gray}{Projet d'étude et de recherche -- UPEC}}
\rfoot{\small\thepage}
\lfoot{\small\textcolor{gray}{Encadrante : Sophie Laruelle}}

%==============================================================
%  MACROS
%==============================================================
\newcommand{\Esp}{\mathbb{E}}
\newcommand{\Var}{\mathrm{Var}}
\newcommand{\Cov}{\mathrm{Cov}}
\newcommand{\xbar}{\bar{x}}
\newcommand{\ybar}{\bar{y}}

%==============================================================
%  PROJET 
%==============================================================
\begin{document}
\renewcommand{\contentsname}{Table des matières}
\renewcommand{\refname}{Références}



%  PAGE DE TITRE

\begin{titlepage}
  \centering
  \vspace*{1.5cm}

  {\Large Université Paris-Est Créteil (UPEC)\\
  Département de Mathématiques\par}

  \vspace{0.8cm}
  \rule{12cm}{0.8pt}
  \vspace{0.8cm}

  {\LARGE\bfseries\color{deepblue}
    Introduction à l'Économétrie\\[0.4em]
    La Régression Linéaire Simple\par}

  \vspace{0.5cm}
  \rule{12cm}{0.8pt}
  \vspace{1cm}

  {\large Projet d'étude et de recherche\par}
  \vspace{0.8cm}

  % Étudiants
  \begin{tcolorbox}[colback=white, colframe=deepblue!60, width=9cm]
    \centering
    {\normalsize\bfseries Étudiants}\\\vspace{0.4em}
    {\normalsize Amine \textsc{El Mir}}\\[0.2em]
    {\normalsize Khalil \textsc{Benaissa}}\\[0.2em]
    {\normalsize Dylan \textsc{Au Yeung}}
  \end{tcolorbox}

  \vspace{0.6cm}
  {\large Encadrante : \textbf{Sophie Laruelle}\par}

  \vspace{1cm}
  \begin{tcolorbox}[colback=deepblue!5, colframe=deepblue, width=12cm,
                    title={\small\bfseries Référence bibliographique principale}]
    \small Gilbert Saporta, \textit{Probabilités, analyse des données et Statistique},
    Éditions TECHNIP, 2006, Chapitre 16.
  \end{tcolorbox}

  \vfill
\end{titlepage}




%  ABSTRACT 

\begin{abstract}
\noindent
Econometrics is the discipline that uses statistical methods to test and quantify
economic theories from empirical observations.
This report is dedicated to the study of simple linear regression, the most fundamental
econometric model, in which a response variable $Y$ is modelled as a linear function
of a single explanatory variable $X$ up to a random noise term~$\varepsilon$.
We first present the theoretical
regression model, distinguishing the probabilistic setting (regression between random
variables) from the fixed-design setting (linear model).
We then derive the Ordinary Least Squares (OLS) estimators $a$ and $b$ of the intercept
$\alpha$ and slope~$\beta$, show that the regression line passes through the empirical
centre of gravity, and establish, via three classical theorems, that these estimators
are unbiased (Theorem~1), of minimum variance among all linear unbiased estimators
(Gauss--Markov theorem, Theorem~2), and that the residual variance estimator $\hat{\sigma}^2$
is also unbiased (Theorem~3).
Statistical tests (analysis of variance, Student test) and a numerical illustration
on a real dataset are also presented.
\end{abstract}

\vspace{0.5cm}
\tableofcontents
\newpage




%  SECTION 1 — INTRODUCTION

\section{Introduction}

L'économétrie est une branche des statistiques appliquées dont l'objet est de mettre
à l'épreuve des théories économiques en ajustant des modèles mathématiques à des
observations. L'exemple emblématique, qui guidera tout ce rapport, est la loi
de Keynes : la consommation $Y$ d'un ménage est supposée être une fonction linéaire
de son revenu $X$. On cherche à tester cette hypothèse et à estimer les paramètres
du modèle sur des données réelles.

Le modèle de régression linéaire simple est le modèle économétrique le plus
élémentaire. Il suppose que l'espérance conditionnelle de $Y$ sachant $X$ est
une fonction affine de $X$ :
\[
  \Esp(Y \mid X) = \alpha + \beta X.
\]
L'objectif est d'estimer les paramètres $\alpha$ (ordonnée à l'origine) et $\beta$
(pente, ou propension marginale à consommer) à partir d'un échantillon
$(x_i, y_i)_{i=1}^n$, ainsi que la variance $\sigma^2$ du terme d'erreur.

Ce rapport est organisé comme suit.
La section~\ref{sec:modele} présente le modèle théorique selon les deux angles
classiques : le cas probabiliste (régression entre variables aléatoires) et le
modèle linéaire à régresseurs fixes.
La section~\ref{sec:mco} dérive les estimateurs des moindres carrés $a$ et $b$.
La section~\ref{sec:proprietes} établit leurs propriétés fondamentales via les
trois théorèmes fondamentaux que nous établirons dans ce rapport.
La section~\ref{sec:tests} aborde les tests statistiques associés.
Enfin, la section~\ref{sec:numerique} illustre la méthode sur un exemple numérique,
avec une implémentation Python.


%  SECTION 2 — MODÈLE THÉORIQUE


\section{Le Modèle Théorique de la Régression Linéaire Simple}

\label{sec:modele}

Dans cette section, nous posons le cadre théorique de la régression linéaire simple
et les hypothèses probabilistes sur lesquelles repose la méthode.
\midskip
\subsection{Le cas probabiliste : régression entre variables aléatoires}

Soient $X$ et $Y$ deux variables aléatoires réelles de carré intégrable.
On sait que l'espérance conditionnelle $\Esp(Y \mid X)$ est la meilleure
approximation de $Y$ par une fonction de $X$ au sens de l'erreur quadratique
moyenne, et que le rapport de corrélation
\[
  \eta^2_{Y/X} = \frac{\Var(\Esp(Y \mid X))}{\Var(Y)}
\]
mesure la qualité de cette approximation.

\begin{definition}[Régression linéaire simple]
\label{def:reglin}
On dit que $Y$ admet une régression linéaire sur $X$ si l'espérance
conditionnelle de $Y$ sachant $X$ est une fonction affine de $X$ :
\begin{equation}
  \label{eq:modeleproba}
  \Esp(Y \mid X) = \alpha + \beta X,
\end{equation}
où $\alpha, \beta \in \mathbb{R}$.
\end{definition}

\noindent On peut alors écrire :
\begin{equation}
  \label{eq:decompY}
  Y = \alpha + \beta X + \varepsilon,
\end{equation}
où $\varepsilon = Y - \Esp(Y \mid X)$ est le résidu aléatoire, vérifiant
$\Esp(\varepsilon) = 0$ et $\Cov(\varepsilon, X) = 0$.
On montre aussi que :
\[
  \beta = \frac{\Cov(X,Y)}{\Var(X)} = \rho \frac{\sigma_Y}{\sigma_X},
  \qquad
  \alpha = \Esp(Y) - \beta\,\Esp(X),
\]
et que dans ce cas $\rho^2 = \eta^2_{Y/X}$, c'est-à-dire que le coefficient de
corrélation linéaire capture entièrement la liaison entre $X$ et $Y$.

\subsection{Le modèle linéaire à régresseurs fixes}


La méthode des moindres carrés s'applique également au cas où les $x_i$ sont des
valeurs fixées non aléatoires, par exemple des périodes d'observations. On parle alors de modèle linéaire plutôt que de régression.

\begin{definition}[Modèle linéaire simple]
\label{def:modlin}
Pour $n$ observations, le modèle linéaire simple s'écrit :
\begin{equation}
  \label{eq:modlin}
  y_i = \alpha + \beta x_i + \varepsilon_i, \quad i = 1, \ldots, n,
\end{equation}
où les $x_i$ sont des constantes données, et les $\varepsilon_i$ sont des
réalisations indépendantes d'une variable d'espérance nulle et de variance
constante $\sigma^2$.
\end{definition}

\noindent Le bon usage de l'estimateur MCO repose sur un ensemble d'hypothèses
portant sur les perturbations $\varepsilon_i$, que nous adoptons dans la suite
en suivant la présentation du chapitre~4 du polycopié de cours~\cite{poly2024} :

\begin{description}
  \item[H1] $\forall i,\; \Esp[\varepsilon_i] = 0$ : la perturbation est centrée.
  \item[H2] $\forall i,\; x_i$ est observé sans erreur (régresseurs fixes).
  \item[H3] $\forall i,\; \Cov(x_i, \varepsilon_i) = 0$ : $X$ et $\varepsilon$ sont
            linéairement indépendants.
  \item[H4] $\forall i,\; \Var(\varepsilon_i) = \sigma^2$ : homoscédasticité des
            perturbations.
  \item[H5] $\forall i \neq j,\; \Cov(\varepsilon_i, \varepsilon_j) = 0$ : les
            perturbations sont non corrélées.
  \item[H6] $\forall i,\; \varepsilon_i \sim \mathcal{N}(0, \sigma^2)$ : normalité
            des perturbations (nécessaire pour les tests exacts).
\end{description}

\begin{remarque}
Les hypothèses H1 à H5 suffisent pour établir le sans-biais et le théorème de
Gauss-Markov. L'hypothèse H6 est requise en plus pour construire des tests
statistiques exacts à partir des lois de Student et de Fisher.
Les propriétés de la méthode des moindres carrés ne dépendent que des lois
conditionnelles à $X$ fixé. On prendra garde de parler de corrélation entre
$Y$ et $X$ que lorsque $X$ est aléatoire.
\end{remarque}




%  SECTION 3 — ESTIMATEURS MCO

\section{Estimation par la Méthode des Moindres Carrés}
\label{sec:mco}

Cette section est dédiée à la construction des estimateurs MCO de $\alpha$ et $\beta$
par minimisation de la somme des carrés des résidus, selon la méthode de Gauss.

\subsection{Principe et critère de minimisation}

On dispose de $n$ couples $(x_i, y_i)$.
La méthode des moindres carrés reprend sur l'échantillon la propriété que
$\Esp(Y \mid X) = \alpha + \beta X$ est la meilleure approximation de $Y$ par $X$
en moyenne quadratique.
On cherche donc à ajuster au nuage de points $(x_i, y_i)$ une droite $y^* = a + bx$
de telle sorte que
\[
  \sum_{i=1}^n (y_i - y_i^*)^2 = \sum_{i=1}^n (y_i - a - bx_i)^2 = F(a,b)
\]
soit minimal.

\subsection{Calcul des estimateurs}

\begin{proposition}[Estimateurs MCO]
\label{prop:mco}
Le minimum de $F(a,b)$ est atteint pour les valeurs $a$ et $b$ solutions du
système des équations normales :
\[
  \frac{\partial F}{\partial a} = 0 \quad \text{et} \quad \frac{\partial F}{\partial b} = 0,
\]
ce qui donne :
\begin{align}
  \label{eq:en1}
  &\sum_{i=1}^n (y_i - a - bx_i) = 0 \quad \Longleftrightarrow \quad \ybar = a + b\xbar, \\
  \label{eq:en2}
  &\sum_{i=1}^n x_i(y_i - a - bx_i) = 0.
\end{align}
La solution est :
\begin{equation}
  \label{eq:b}
  {b = \frac{\displaystyle\sum_{i=1}^n (x_i - \xbar)(y_i - \ybar)}
               {\displaystyle\sum_{i=1}^n (x_i - \xbar)^2}
            = r \frac{s_y}{s_x}},
  \qquad
  {a = \ybar - b\xbar},
\end{equation}
où $r$ est le coefficient de corrélation linéaire empirique, $s_x$ et $s_y$ sont
les écarts-types empiriques de $X$ et $Y$.
\end{proposition}
\smallskip
\begin{demonstration} 
On applique les hypothèses :
\[
  \frac{\partial F}{\partial a} = -2\sum_{i=1}^n(y_i - a - bx_i) = 0
  \quad \Rightarrow \quad \sum_i y_i = na + b\sum_i x_i,
  \quad \text{soit} \quad \ybar = a + b\xbar.
\]
\[
  \frac{\partial F}{\partial b} = -2\sum_{i=1}^n x_i(y_i - a - bx_i) = 0
  \quad \Rightarrow \quad \sum_i x_i y_i = a\sum_i x_i + b\sum_i x_i^2.
\]
Avec $a = \ybar - b\xbar$ dans la seconde équation on a :
\[
\sum_i x_i y_i = (\ybar - b\xbar) \sum_i x_i + b\sum_i x_i^2 \quad \Rightarrow \quad \sum_i x_i y_i -n\xbar\ybar = -b\xbar^2n + b \sum_i x_i^2
\]
\[
\ \Rightarrow \quad
  \sum_i x_i y_i - n\xbar\ybar = b\!\left(\sum_i x_i^2 - n\xbar^2\right),
  \quad \text{ainsi} \quad
  \sum_i (x_i - \xbar)(y_i - \ybar) = b\sum_i(x_i-\xbar)^2,
\]

d'où la formule~\eqref{eq:b}. 
La fonction $F$ est strictement convexe en $(a,b)$
donc ce point critique est bien un minimum global.
\end{demonstration}

\begin{remarque}[Propriété géométrique]
L'équation~\eqref{eq:en1} montre que la droite des moindres carrés passe par le
centre de gravité du nuage $(\xbar, \ybar)$. Qui s'écrit :
\[
  y^* = \ybar + b(x - \xbar).
\]
\end{remarque}




%  SECTION 4 — PROPRIÉTÉS

\section{Propriétés des Estimateurs MCO}
\label{sec:proprietes}

Nous établissons dans cette section les trois théorèmes fondamentaux
de l'estimation par moindres carrés ordinaires. On note $A$ et $B$ les variables
aléatoires dont $a$ et $b$ sont les réalisations.

Notons d'abord que $b$ s'écrit plus simplement :
\begin{equation}
  \label{eq:b_simplifie}
  b = \frac{\sum_{i=1}^n (x_i - \xbar) y_i}{\sum_{i=1}^n (x_i - \xbar)^2},
\end{equation}
qui résulte du fait que $\sum_i (x_i - \xbar) = 0$.

  
\subsection{Théorème 1 (Sans biais)}

\begin{theoreme}[Sans biais]
\label{thm:unbiased}
Sous le modèle~\eqref{eq:modlin}, $a$, $b$ et $y^* = a + bx$ sont des estimations
sans biais de $\alpha$, $\beta$ et de $\Esp(Y \mid X = x) = \alpha + \beta x$ :
\[
  \Esp(A) = \alpha, \qquad \Esp(B) = \beta, \qquad \Esp(Y^*) = \alpha + \beta x.
\]
\end{theoreme}

\begin{demonstration}
Nous montrons que $\Esp_{(x)}(B) = \beta$, où $\Esp_{(x)}$ désigne l'espérance
conditionnelle sachant les $x_i$.
D'après~\eqref{eq:b_simplifie} et la linéarité de l'espérance :
\begin{align*}
  \Esp_{(x)}(B)
  &= \frac{\sum_{i=1}^n (x_i - \xbar)\,\Esp_{(x)}(Y_i)}{\sum_{i=1}^n (x_i - \xbar)^2}
   = \frac{\sum_{i=1}^n (x_i - \xbar)(\alpha + \beta x_i)}{\sum_{i=1}^n (x_i - \xbar)^2}.
\end{align*}
Or $\sum_i (x_i - \xbar)\cdot\alpha = \alpha\sum_i(x_i-\xbar) = 0$, donc :
\[
  \Esp_{(x)}(B) = \frac{\beta\sum_{i=1}^n (x_i - \xbar) x_i}
                       {\sum_{i=1}^n (x_i - \xbar)^2}
               = \frac{\beta\sum_{i=1}^n (x_i - \xbar)^2}
                       {\sum_{i=1}^n (x_i - \xbar)^2} = \beta.
\]
L'égalité $\Esp(B) = \Esp(\Esp_{(x)}(B)) = \beta$ en découle.
Comme $A = \ybar - BX$, on a :
\[
  \Esp_{(x)}(A) = \Esp_{(x)}(\bar{Y}) - \xbar\,\Esp_{(x)}(B)
                = (\alpha + \beta\xbar) - \xbar\beta = \alpha.
\]
Donc $\Esp(A) = \alpha$, et $\Esp(Y^*) = \Esp(A) + \Esp(B)\cdot x = \alpha + \beta x$.
\end{demonstration}

\begin{remarque}
Le fait d'être sans biais n'est qu'une qualité mineure pour des estimateurs.
Le théorème suivant établit une propriété d'optimalité bien plus forte.
\end{remarque}

\subsection{Variances des estimateurs}

Avant d'énoncer le théorème de Gauss-Markov, calculons les variances de $A$ et $B$.

\begin{proposition}[Variances de $A$ et $B$]
\label{prop:variances}
Sous le modèle~\eqref{eq:modlin} avec $\Var(\varepsilon_i) = \sigma^2$
et $\Cov(\varepsilon_i, \varepsilon_j) = 0$ pour $i \neq j$ :
\begin{equation}
  \label{eq:varB}
  \Var(B) = \frac{\sigma^2}{\displaystyle\sum_{i=1}^n (x_i - \xbar)^2},
\end{equation}
\begin{equation}
  \label{eq:varA}
  \Var(A) = \sigma^2\!\left(\frac{1}{n} + \frac{\xbar^2}{\displaystyle\sum_{i=1}^n (x_i - \xbar)^2}\right).
\end{equation}
\end{proposition}

\begin{demonstration}
En notant $c_i = (x_i - \xbar)/\sum_j (x_j - \xbar)^2$, on a $B = \sum_i c_i Y_i$,
donc par indépendance des $\varepsilon_i$ :
\[
  \Var(B) = \sigma^2 \sum_{i=1}^n c_i^2
          = \frac{\sigma^2\,\sum_i(x_i-\xbar)^2}{\left(\sum_i(x_i-\xbar)^2\right)^2}
          = \frac{\sigma^2}{\sum_i(x_i-\xbar)^2}.
\]
Puis, comme $A = \bar{Y} - B\xbar$ et $B$, $\bar{Y}$ sont non corrélés
:
\[
  \Var(A) = \Var(\bar{Y}) + \xbar^2\Var(B) = \frac{\sigma^2}{n} + \frac{\sigma^2\xbar^2}{\sum_i(x_i-\xbar)^2}.
\]
\end{demonstration}

\begin{remarque}[Covariance entre $A$ et $B$]
On montre de même que~\cite{poly2024} :
\[
  \Cov(A, B) = -\xbar\,\Var(B) = \frac{-\xbar\,\sigma^2}{\displaystyle\sum_{i=1}^n(x_i - \xbar)^2}.
\]
Ainsi, lorsque $\xbar > 0$, la pente et l'ordonnée à l'origine estimées sont
corrélées négativement : si $b$ augmente, $a$ tend à diminuer (et réciproquement),
ce qui traduit la contrainte que la droite passe par le centre de gravité $(\xbar, \ybar)$.
\end{remarque}

\subsection{Optimalité : le théorème de Gauss-Markov (Théorème 2)}

\begin{theoreme}[Gauss-Markov]
\label{thm:gm}
Sous le modèle~\eqref{eq:modlin}, parmi tous les estimateurs sans biais de $\alpha$
et $\beta$ qui sont des fonctions linéaires des $Y_i$, les estimateurs MCO $A$ et $B$
sont ceux de variance minimale : ce sont les MSLB (Meilleurs estimateurs Sans biais Linéaires).
\end{theoreme}

\begin{demonstration}
Nous montrons le résultat pour $B$ ; le cas de $A$ est analogue.
Soit $\widetilde{B} = \sum_i \ell_i Y_i$ un estimateur linéaire sans biais quelconque.
La condition de sans-biais donne :
\[
  \Esp_{(x)}(\widetilde{B}) = \sum_i \ell_i(\alpha + \beta x_i) = \beta
  \quad \forall \alpha, \beta
  \quad \Longleftrightarrow \quad
  \sum_i \ell_i = 0 \quad \text{et} \quad \sum_i \ell_i x_i = 1.
\]
Posons $d_i = \ell_i - c_i$ où $c_i = (x_i - \xbar)/\sum_j(x_j-\xbar)^2$.
Les conditions de sans-biais deviennent $\sum_i d_i = 0$ et $\sum_i d_i x_i = 0$.
On calcule :
\[
  \Var(\widetilde{B}) = \sigma^2\sum_i(\ell_i)^2 = \sigma^2\sum_i(c_i + d_i)^2
  = \sigma^2\!\left(\sum_i c_i^2 + 2\sum_i c_i d_i + \sum_i d_i^2\right).
\]
Le terme croisé vaut :
\[
  \sum_i c_i d_i = \frac{\sum_i(x_i-\xbar)d_i}{\sum_j(x_j-\xbar)^2}
                = \frac{\sum_i x_i d_i - \xbar\sum_i d_i}{\sum_j(x_j-\xbar)^2} = 0
\]
par les conditions de sans-biais.
Donc :
\[
  \Var(\widetilde{B}) = \sigma^2\sum_i c_i^2 + \sigma^2\sum_i d_i^2
                      = \Var(B) + \sigma^2\sum_i d_i^2 \geq \Var(B),
\]
avec égalité si et seulement si $d_i = 0$ pour tout $i$, c'est-à-dire
$\widetilde{B} = B$.
\end{demonstration}




%  SECTION 5 — ESTIMATION DE SIGMA^2

\section{Estimation de la Variance des Erreurs}
\label{sec:sigma}

Cette section présente la construction et le sans-biais de l'estimateur de $\sigma^2$.

\begin{definition}[Résidus et variance résiduelle]
Les écarts résiduels sont $e_i = y_i - y_i^* = y_i - a - bx_i$.
La variance résiduelle empirique est $s_{y/x}^2 = \frac{1}{n}\sum_{i=1}^n e_i^2$.
\end{definition}

On établit d'abord deux propriétés des résidus.

\begin{proposition}[Propriétés des résidus]
\label{prop:resid}
Les résidus $e_i$ vérifient : $\sum_{i=1}^n e_i = 0$ et $\sum_{i=1}^n x_i e_i = 0$.
\end{proposition}

\begin{demonstration}
Ces deux identités sont exactement les équations normales~\eqref{eq:en1}--\eqref{eq:en2},
réécrites en termes des résidus $e_i = y_i - a - bx_i$.
\end{demonstration}

\begin{remarque}
La nullité de $\sum e_i$ montre que les résidus ne sont pas des réalisations
indépendantes : ils sont soumis à deux contraintes linéaires.
\end{remarque}

\begin{theoreme}[Sans biais de $\hat{\sigma}^2$]
\label{thm:sigma}
L'estimateur
\begin{equation}
  \label{eq:shat}
  \hat{\sigma}^2 = \frac{\displaystyle\sum_{i=1}^n e_i^2}{n - 2}
                = \frac{n\,s_{y/x}^2}{n-2}
\end{equation}
est un estimateur sans biais de $\sigma^2$ :
$\Esp[\hat{\sigma}^2] = \sigma^2$.
\end{theoreme}

\begin{demonstration}
Exprimons $e_i$ en fonction des erreurs vraies $\varepsilon_i$ :
\[
  e_i = y_i - a - bx_i = (\alpha + \beta x_i + \varepsilon_i)
        - (\ybar - b\xbar) - b x_i
      = \varepsilon_i - \bar{\varepsilon} - (b - \beta)(x_i - \xbar),
\]
où $\bar{\varepsilon} = \frac{1}{n}\sum_i \varepsilon_i$.
En élevant au carré et en sommant :
\[
  \sum_i e_i^2
  = \sum_i\varepsilon_i^2 - n\bar{\varepsilon}^2 - (b-\beta)^2\sum_i(x_i-\xbar)^2.
\]
En prenant l'espérance et en utilisant
$\Esp[\bar{\varepsilon}^2] = \sigma^2/n$ et
$\Esp[(b-\beta)^2] = \Var(B) = \sigma^2/\sum_i(x_i-\xbar)^2$ :
\[
  \Esp\!\left[\sum_i e_i^2\right]
  = n\sigma^2 - \sigma^2 - \sigma^2 = (n-2)\sigma^2.
\]
Donc $\Esp[\hat{\sigma}^2] = \sigma^2$.
\end{demonstration}

\begin{remarque}
Le facteur $n-2$ traduit la perte de deux degrés de liberté
consécutive à l'estimation de deux paramètres $\alpha$ et $\beta$.
\end{remarque}




%  SECTION 6 — TESTS STATISTIQUES

\section{Tests Statistiques dans le Modèle Linéaire}
\label{sec:tests}

Sous l'hypothèse supplémentaire de normalité des erreurs,
$\varepsilon_i \sim \mathcal{N}(0, \sigma^2)$, on dispose de tests exacts.

\subsection{Analyse de variance de la régression}

On dispose de la décomposition de la variance:
\begin{equation}
  \label{eq:decomp}
  \underbrace{\sum_{i=1}^n (y_i - \ybar)^2}_{\text{SCT (totale)}}
  =
  \underbrace{\sum_{i=1}^n (y_i^* - \ybar)^2}_{\text{SCE (expliquée)}}
  +
  \underbrace{\sum_{i=1}^n e_i^2}_{\text{SCR (résiduelle)}}.
\end{equation}
Cette décomposition repose sur l'orthogonalité $\sum_i(y_i^*-\ybar)(y_i - y_i^*) = 0$.
On définit le coefficient de détermination $R^2 = \text{SCE}/\text{SCT} = r^2$.

Sous l'hypothèse $H_0 : \beta = 0$ (non-régression), on montre que~\cite{poly2024} :
\[
  F = \frac{\text{SCE}/1}{\text{SCR}/(n-2)}
    = \frac{(n-2)\,R^2}{1 - R^2}
    \sim \mathcal{F}(1;\, n-2).
\]
La seconde expression, exprimée en fonction de $R^2$, montre directement que
rejeter $H_0$ revient à tester si le coefficient de détermination est
significativement non nul.
Le tableau d'analyse de variance récapitule ces résultats :

\begin{center}
\begin{tabular}{lccc}
\toprule
\textbf{Source} & \textbf{SC} & \textbf{ddl} & \textbf{Carré moyen} \\
\midrule
Régression & SCE & $1$ & SCE \\
Résiduelle & SCR & $n-2$ & SCR$/(n-2)$ \\
\midrule
Totale & SCT & $n-1$ & \\
\bottomrule
\end{tabular}
\end{center}

\subsection{Test de Student sur $\beta$}

Un test équivalent consiste à tester $H_0 : \beta = 0$ contre $H_1 : \beta \neq 0$
via la statistique :
\[
  t = \frac{b}{\hat{\sigma}/\sqrt{\sum_i(x_i-\xbar)^2}} \sim \mathcal{T}(n-2)
  \quad \text{sous } H_0.
\]
Ce test est strictement équivalent au test $F$ ci-dessus et au test $H_0 : \rho = 0$
sur le coefficient de corrélation~\cite{saporta2006}.




%  SECTION 7 — ILLUSTRATION NUMÉRIQUE

\section{Illustration Numérique}
\label{sec:numerique}


\subsection{Modèle d'Evaluation des Actifs Financiers appliquée à LVMH}

Le MEDAF est un modèle fondamental de la finance de marché, il établit que la prime de risque individuelle d'un actif, c'est à dire l'excès de rendement $E(R_i) - R_f$ par rapport au taux sans risque $R_f$  est proportionnelle à la prime de risque du marché $E(R_m) - R_f$ :
\begin{equation}
    E(R_i)= R_f + \beta_i(E(R_m)-R_f)
\end{equation}
le coefficient $\beta_i$ mesure la sensibilité de l'actif aux fluctuations du marché. Le terme $\alpha_i$ est l'ordonnée à l'origine de la droite de régression, lorsque ce terme est strictement positif il indique un actif sous-évalué, si il est strictement négatif un actif surévalué et si il est nulle alors une valorisation juste.
On peut alors montrer que \begin{equation}
    \beta_i = \rho_{i,m} \frac{\sigma_i}{\sigma_m}
\end{equation}
où $\rho_{i,m}$ désigne le coefficient de corrélation entre l'actif et le marche, $\sigma_i$ et  $\sigma_m$ sont leurs écart types respectifs on retrouve la formule de l'estimateur MCO.

Nous estimons ce modèle sur $n = 48$ rendements hebdomadaires de l'action
LVMH et de l'indice CAC40, sur la période janvier 2024 -- décembre 2024.
Les rendements sont calculés à partir des prix de clôture $P_t$ par :
\[
    r_t = \frac{P_t - P_{t-1}}{P_{t-1}}, \quad t = 1, \ldots, 51.
\]

\begin{center}
\begin{tabular}{ccccc}
\toprule
$t$ & $P_t^{\text{LVMH}}$ en euros & $P_t^{\text{CAC}}$ & 
$r_t^{\text{LVMH}}$ (en \%) & $r_t^{\text{CAC}}$ ( en \%) \\
\midrule
0 & 669,90 & 7465,14 & — & — \\
1 & 659,40 & 7371,64 & $-1{,}57$ & $-1{,}25$ \\
2 & 773,10 & 7634,14 & $17{,}24$ & $3{,}56$ \\
3 & 778,10 & 7592,26 & $0{,}65$ & $-0{,}55$ \\
4 & 807,50 & 7647,52 & $3{,}78$ & $0{,}73$ \\
5 & 815,80 & 7768,18 & $1{,}03$ & $1{,}58$ \\
\bottomrule
\end{tabular}
\end{center}

Les statistiques descriptives sur les $n = 48$ rendements sont :
\[
\bar{x} = 0,0038\%, \quad \bar{y} = 0,0129\%, \quad
s_x = 1,9622\%, \quad s_y = 4,8450\%, \quad r = 0,7436.
\]

On en déduit les estimateurs MCO :
\[
\hat{\beta} = r \cdot \frac{s_y}{s_x} = 0,7436 \times \frac{4,8450}{1,9622} = 1,8361,
\qquad
\hat{\alpha} = \bar{y} - \hat{\beta}\,\bar{x} = 0,0059\%.
\]
Le coefficient de détermination vaut $R^2 = 0,5529$ et l'écart-type résiduel
$\hat{\sigma} = 3,31\%$. On teste l'hypothèse nulle $H_0 : \beta = 0$  contre $H_1 : \beta \neq 0$. La statistique 
de Student donne :
\[
    t = \frac{\hat{\beta}}{\widehat{SE}(\hat{\beta})} =
    \frac{1,8361}{0,2434} = 7,54 > t_{97,5\%}^{(46)} = 2,01
\]
avec une p-value de $1,41 \times 10^{-9}$ : on rejette $H_0 : \beta = 0$
au seuil $5\%$. L'intervalle de confiance à $95\%$ pour $\beta$ est
$[1,346\ ;\ 2,326]$.

\begin{remarque}
Le bêta estimé $\hat{\beta} = 1{,}84 > 1$ indique que LVMH est plus
volatile que le marché sur cette période : une hausse de $1\%$ du CAC40
est associée en moyenne à une hausse de $1,84\%$ de LVMH. Le $R^2 = 0{,}55$
signifie que le CAC40 explique $55\%$ de la variabilité des rendements LVMH,
le reste étant dû à des facteurs propres à l'entreprise. Enfin,
$\hat{\alpha} \approx 0$ est cohérent avec l'hypothèse du MEDAF qui prédit
$\alpha = 0$ pour un actif correctement valorisé.
\end{remarque}

\subsection{Loi de Keynes sur données de l'INSEE}

Pour illustrer d'avantage la méthode et faire le lien avec la problématique de l'introduction, nous appliquons la régression linéaire aux données de l'INSEE
sur l'évolution des dépenses de consommation finale ($Y$) et du revenu brut ($X$) des ménages français entre 2010 et 2019~\cite{poly2024}. Les valeurs
sont exprimées en pourcentage de variation par rapport à l'année précédente.

\begin{center}
\begin{tabular}{ccc}
\toprule
\textbf{Année} & \textbf{Consommation $Y$ (\%)} & \textbf{Revenu $X$ (\%)} \\
\midrule
2010 & 3,0 & 2,8 \\
2011 & 2,4 & 2,0 \\
2012 & 0,9 & 1,0 \\
2013 & 1,1 & $-0,5$ \\
2014 & 0,9 & 1,3 \\
2015 & 1,7 & 1,1 \\
2016 & 2,0 & 1,8 \\
2017 & 2,3 & 2,6 \\
2018 & 2,6 & 3,0 \\
2019 & 2,5 & 3,1 \\
\bottomrule
\end{tabular}
\end{center}

Les statistiques descriptives donnent $n = 10$, $\xbar = 1{,}82$, $\ybar = 1{,}94$,
$s_x^2 = 1{,}1476$, $s_y^2 = 0{,}5144$ et $s_{xy} = 0{,}6412$.
Les estimateurs MCO sont :
\[
  \hat{b} = \frac{s_{xy}}{s_x^2} = \frac{0{,}6412}{1{,}1476} \approx 0{,}559
  \qquad
  \hat{a} = \ybar - \hat{b}\,\xbar = 1{,}94 - 0{,}559 \times 1{,}82 \approx 0{,}923
\]
Le coefficient de détermination vaut $R^2 \approx 0{,}697$, soit un coefficient de
corrélation $r \approx 0{,}835$.
Le test de Student donne $t \approx 4{,}28 > t_{0{,}975}^{(8)} = 2{,}306$ donc rejet de 
 $H_0 : \beta = 0$ au seuil $5\%$.

\begin{remarque}
Ce résultat confirme la loi de Keynes évoquée en introduction : la consommation
croît avec le revenu ($\hat{b} > 0$), mais avec une propension marginale à
consommer inférieure à~1 ($\hat{b} \approx 0{,}56 < 1$). Le $R^2 = 0{,}70$
indique que le revenu explique environ $70\%$ de la variabilité de la consommation
sur cette période, ce qui est cohérent.
\end{remarque}



%  SECTION 8 — CONCLUSION

\section{Conclusion}

Nous avons étudié le modèle de régression linéaire simple $y_i = \alpha + \beta x_i + \varepsilon_i$ et ses propriétés fondamentales.

Les estimateurs MCO $a$ et $b$ s'obtiennent par minimisation de la somme des carrés des
résidus ; la droite ajustée passe par le centre de gravité $(\xbar, \ybar)$ du nuage
et sa pente vaut $b = r\,s_y/s_x$. Les trois théorèmes centraux garantissent que :
\begin{itemize}
  \item $a$, $b$ et $y^*$ sont \textbf{sans biais} (Théorème~1) ;
  \item $A$ et $B$ sont les \textbf{estimateurs de variance minimale} parmi tous les
        estimateurs linéaires sans biais --- théorème de Gauss-Markov (Théorème~2) ;
  \item $\hat{\sigma}^2 = \sum e_i^2/(n-2)$ est un estimateur \textbf{sans biais} de
        $\sigma^2$ (Théorème~3), le dénominateur $n-2$ reflétant la perte de deux
        degrés de liberté.
\end{itemize}
Sous l'hypothèse de normalité des erreurs, la statistique de Student $t = b/\hat{\sigma}_b$
suit une loi $\mathcal{T}(n-2)$ sous $H_0 : \beta = 0$, ce test étant équivalent au
test $F$ d'analyse de variance et au test sur le coefficient de corrélation.




%  BIBLIOGRAPHIE

\begin{thebibliography}{9}

\bibitem{saporta2006}
G.~Saporta,
\textit{Probabilités, analyse des données et Statistique},
Éditions TECHNIP, Paris, 2006.
Chapitre~16 : La régression simple, page~387--405.

\bibitem{wooldridge2019}
J.~M.~Wooldridge,
\textit{Introductory Econometrics: A Modern Approach},
7\textsuperscript{e} édition, Cengage Learning, 2019.

\bibitem{casella2002}
G.~Casella et R.~L.~Berger,
\textit{Statistical Inference},
2\textsuperscript{e} édition, Duxbury Press, 2002.

\bibitem{poly2024}
\textsc{S. Laruelle},
\textit{Introduction à l'Économétrie},
Polycopié de cours, Chapitre~4, Université Paris-Est Créteil (UPEC).

\end{thebibliography}

\end{document}